\section{Restoring the sparse argument without a logarithmic lower bound}
\label{sec:sparse-restored}

This revision restores a proof chain that was present in the 9 October
degree-plus-spectrum draft\footnote{The retained source is
\href{https://github.com/zhixin0825/spectral-lrt/blob/de60b64c7d7f127ae6dc1d010c840af707a080da/paper/previous_draft/sbm_spectral_lrt_20261009.tex}{the 9 October sparse proof at commit \texttt{de60b64}}.}:
a same-threshold truncated-row-energy
recursion and local full-to-deleted-graph comparison. That draft treated
$d\to\infty$, $d=o(\log n)$ and obtained a leading-exponent result.
Its non-tail score errors were $O(\sqrt d)$ or $o(d)$, so it cannot be
quoted verbatim for a constant risk ratio. The argument below reuses its
row-tail mechanism, proves a scalar $O(1)$ LOO comparison for the current
floor normalization, and supplies the fitted-parameter bounds.

The distinction between decoders is explicit. The completed sparse
result below uses the finite radius-count initializer already specified
for DC-SBM and block-parameter Bernoulli likelihood updates. Its ordinary
SBM specialization fixes all degree parameters at one and uses the
ordinary, unprofiled Bernoulli LRT. It is not a theorem for the unchanged
ordinary free-profile growing/replacement/soft-EM schedule of
\cref{sbm:algorithm}. That schedule's global generated-profile box must
not be transferred silently to graphs with exceptional high degrees.

\subsection{The regression caused by whole-graph fallback}

\begin{proposition}[A whole-graph zero-degree fallback is invalid in the
sparse regime]\label{restore:fallback}
Consider a balanced two-group SBM with within probability $p$, between
probability $p/2$, and $d=np=\log\log n$. The probability of at least
one isolated vertex tends to one. For every deterministic fallback
partition fixed in advance, there is a balanced true partition for which
an algorithm returning that fallback whenever an isolated vertex exists
has expected permutation-invariant misclassification bounded below by
$1/4+o(1)$. Its local oracle risk tends to zero.
\end{proposition}

\begin{proof}
For every vertex the isolation probability is
$(1-p)^{n/2-1}(1-p/2)^{n/2}
=\exp\{-3d/4+o(1)\}$. If $N_0$ counts isolated vertices, then
$\E N_0=n\exp\{-3d/4+o(1)\}\to\infty$. For distinct vertices,
their two isolation events share only the edge between them, giving
\[
 \frac{\operatorname{Var}(N_0)}{(\E N_0)^2}
 \le \frac1{\E N_0}+\frac p{1-p}\longrightarrow0.
\]
Choose a balanced true partition whose overlap with the fixed fallback
has error at least $1/4$ after either label permutation. Such a partition
exists by splitting each fallback class approximately evenly. The loss
lower bound follows. Standard incident-row Chernoff comparison gives
oracle risk at most $Ce^{-cd}\to0$ on this separated model.
\end{proof}

Thus the logarithmic restriction cannot be removed merely by deleting
text from the unregularized theorem. The valid repair is to restore a
positive floor and local risk accounting. With the regularized encoder,
a zero-degree vertex has a zero reconstructed row and causes no
whole-graph fallback.

\subsection{The regularized finite algorithm}
\label{restore:algorithm}

Fix sufficiently small $a,\alpha>0$ and sufficiently large $\beta>0$,
depending only on the fixed structural parameter bounds. Choose them so
that the population degrees lie strictly inside the propensity-clamping
interval, and the population floor is inactive. Put
\[
 \mathbf d=A\one,\quad \bar D=n^{-1}\one^\top\mathbf d,\quad
 \tau=a\bar D,\quad G_i=\max(D_i,\tau).
\]
The all-zero graph has a fixed output. Otherwise $G$ is positive, even
when some vertices are isolated. Retain the $K$ eigenpairs with largest
absolute eigenvalues of $S_G=G^{-1/2}AG^{-1/2}$ and the raw degrees.
There is no requirement that a Perron degree vector be one of these
eigenvectors after flooring. Compute
\[
 R_0=G^{1/2}U\Lambda U^\top G^{1/2}=AH_G,
 \qquad H_G=G^{-1/2}UU^\top G^{1/2},
\]
and use the exact-degree, hollow reconstruction $R^h$ defined in
\eqref{restore:reconstruction}. Only $(\mathbf d,U,\Lambda)$ are used
after encoding; $G$ is recovered from $\mathbf d$.

Initialize by normalizing the rows of $U$ and using the deterministic
radius-$1/8$, deletion-radius-$1/4$ count procedure in
\cref{dcsbm:initializer}. A zero row is assigned any fixed unit direction
for this initialization, without stopping other vertices.

For DC-SBM, define
\[
 T_i=\min\{\max(D_i,\alpha\bar D),\beta\bar D\},\qquad
 \widehat\theta_i=T_i/\overline T_{\widehat z_i}.
\]
Fit the symmetric block matrix by the exact signed Bernoulli working
criterion \eqref{dcsbm:block-criterion}, with $R^h$ and these degree
estimates, under the same fixed candidate-probability cap. Classify each
row by exact Bernoulli target profiling, with the degree parameter
profiled separately for every candidate community.

For ordinary SBM, set $\widehat\theta_i=1$ and maximize the same block
criterion; use the ordinary Bernoulli row likelihood, without target
profiling. Equivalently each fitted symmetric block probability is its
ordered reconstructed block total divided by the corresponding number
of ordered, self-excluded pairs, clipped to
$[\varepsilon,q_{\max}]$, where
$\varepsilon=c_0\bar D/n$ for a fixed sufficiently small $c_0>0$.
On the proved event the maximizers are interior, so this clipping has
no statistical effect.

At least one fit-and-LRT update is performed. The procedure may use any
finite number of these hard refits, including finite data-dependent
stopping. The invariant basin below covers every such endpoint
simultaneously. An undefined empty-group or parameter-fit operation may
have a fixed fallback; those operations are well-defined on the proved
aggregate event. No individual zero degree triggers a graph-wide fallback.

\subsection{Truncated spectral-row energy for the empirical degree floor}
\label{restore:tail-section}

This subsection restores the truncated-row argument for the actual
degree-floor encoder.  The normalization requires a square-root cap,
rather than the linear cap used in the earlier sparse construction.
We include its additional degree-tail term explicitly.  No uniform
row-incoherence event or logarithmic degree lower bound is imposed.

Work under the fixed structural assumptions of
\cref{prop:sparse-floor}, with $d=np\to\infty$ and $np^2\le C_1$.
Use the empirical version of the floor described after that proposition:
\[
 \tau=a_0\bar D,\qquad
 G_j=\max(D_j,\tau),\qquad
 W=\tau^{1/2}G^{-1/2},\qquad M=WAW=\tau S_G.
\]
The constant $a_0>0$ is fixed and sufficiently small for the structural
class.  At the all-zero graph use the prescribed fixed output.
Let $U$ be an orthonormal basis of the informative $K$-dimensional
eigenspace and write $MU=U\Xi$, so $\Xi=\tau\Lambda$.
All bounds below are uniform in the fixed compact structural class.
In particular,
\begin{equation}
 0\le\Omega_{ij}\le Cd/n,\qquad
 \norm{\Omega_{i\cdot}}_2\le Cd/\sqrt n.
 \label{restore:tail-population-rows}
\end{equation}

Fix a probability parameter $r>0$.  Apply
\cref{prop:sparse-floor} with a sufficiently larger fixed parameter to
the full graph and all graphs obtained by deleting at most two vertices.
Let $\mathcal E_r$ be the intersection of the resulting geometry and
two-sided total-degree events.  Restricting the population matrix
preserves the fixed structural bounds and comparable block masses.
Thus
\begin{equation}
 \PP(\mathcal E_r^c)\le C_rn^{-r}.
 \label{restore:tail-geometry-probability}
\end{equation}
On this event, for every graph in the deletion family, its recomputed
threshold is comparable to $d$, its total degree is at most $C_rnd$,
its centered weighted noise has norm at most $C_r\sqrt d$, its
informative eigenvalues have absolute values at least $c_rd$, and its
remaining eigenvalues have absolute values at most $C_r\sqrt d$.
These are consequences of the aggregate geometry already proved; a
bound on the largest spectral row is not included in
$\mathcal E_r$.  When a particular deleted graph is used as independent
training data, use its own analogous event formed from that graph and
its further deletions.  That event is measurable from the training
graph alone and has the same probability bound.

Two deterministic properties of the floor will be useful:
\begin{equation}
 W_j\le1,\qquad W_j^2D_j\le\tau,
 \qquad \norm M_{\op}\le\tau.
 \label{restore:tail-floor-operator}
\end{equation}
Indeed, $G\ge\diag(A\one)$ and the inequality
$2|ab|\le a^2+b^2$ gives
$|x^\top S_Gx|\le\sum_j(D_j/G_j)x_j^2\le\norm x_2^2$.

\begin{lemma}[Truncated energy for the degree-floor eigenvectors]
\label{restore:tail-energy-lemma}
There are fixed constants $R_0,C,c>0$ such that, for every fixed
$R\ge R_0$ and all sufficiently large $n$,
\begin{equation}
 \E\!\left[\ind{\mathcal E_r}
       \sum_{j=1}^n\norm{U_{j\cdot}}_2^2
                  \ind{\sqrt n\norm{U_{j\cdot}}_2>R}\right]
 \le CR^2e^{-cd\log R}.
 \label{restore:tail-energy}
\end{equation}
The constants $R_0,C,c$ depend on the fixed structural bounds,
$a_0$, and $r$, but not on $R$ or $n$.  The sufficiently-large-$n$
condition may depend on the chosen fixed $R$.  The same result holds
for each recomputed deleted graph with its training-measurable
geometry event.
\end{lemma}

\begin{proof}
We give the weight perturbation, the same-threshold energy recurrence,
and its expectation bounds separately.

\medskip\noindent\textbf{Recomputed weights and the eigenspace residual.}
For $J=[n]\setminus\{i\}$, let $W^{(-i)}$ and $U^{(-i)}$ be
computed from $A_{JJ}$, with $U^{(-i)}$ embedded by a zero row at $i$.
Outside the geometry event choose an arbitrary deterministic
orthonormal fallback depending only on $A_{JJ}$.
For an all-zero training graph set $W^{(-i)}=I$ and use that
orthonormal fallback. Thus the proof-only completion always has
$\norm{W^{(-i)}}_{\op}\le1$, and the independent innovation below
is defined for every training graph.
Put $(\Gamma_i)_{jj}=W_j^{(-i)}/W_j-1$ for $j\in J$ and
$(\Gamma_i)_{ii}=0$.
The exact threshold difference is
\[
 \tau^{(-i)}-\tau
 =a_0\frac{\bar D-2D_i}{n-1}.
\]
Comparing the two square-root caps, using the comparability of the
thresholds to $d$, gives on $\mathcal E_r$
\begin{equation}
 |(\Gamma_i)_{jj}|
 \le C\left\{\frac{A_{ij}}{G_j}
             +\frac{d+D_i}{nd}\right\},
 \qquad \norm{\Gamma_i}_{\op}\le C/d.
 \label{restore:tail-weight-change}
\end{equation}
This also covers a crossing of a floor threshold.  It follows either
by splitting the cases in the two maxima or by comparing the logarithms
of $\sqrt{\tau/\max(D_j,\tau)}$; the relative threshold change is
$O((d+D_i)/(nd))$.  The last bound uses $D_i\le n$ and $d\le Cn$.

There is an orthogonal matrix $O_i$ such that
\begin{equation}
 \norm{U-U^{(-i)}O_i}_F
 \le C\norm{U_{i\cdot}}_2+C\norm{\Gamma_iU}_F.
 \label{restore:tail-loo-residual}
\end{equation}
To verify this, write
$M=\left(\begin{smallmatrix}0&m_i^\top\\m_i&B_i\end{smallmatrix}\right)$.
By \eqref{restore:tail-floor-operator},
$\norm{m_i}_2\le\sqrt\tau$ and $\norm{B_i}_{\op}\le\tau$.
Also $B_iU_J=U_J\Xi-m_iU_{i\cdot}$.
The lower block of the recomputed deleted matrix is
$(I+\Gamma_i)B_i(I+\Gamma_i)$.
Subtracting the full matrix and applying the difference to $U$ gives
a Frobenius residual bounded by
$Cd(\norm{U_{i\cdot}}_2+\norm{\Gamma_iU}_F)$.
The full informative spectrum is separated from the deleted
noninformative spectrum by $cd$.
The complementary-space Sylvester equation and orthogonal Procrustes
alignment give \eqref{restore:tail-loo-residual}.
Only separation from the complementary spectrum is used; the
informative eigenvalues can repeat or have either sign.

\medskip\noindent\textbf{Truncation and summable weight changes.}
Set $t=R/\sqrt n$.  Let $f_t$ be radial projection onto the Euclidean
ball of radius $t$, applied to each row, and put
\[
 V=f_t(U),\qquad T=U-V,\qquad
 H_R(U)=\sum_j\norm{U_{j\cdot}}_2^2
                         \ind{\norm{U_{j\cdot}}_2>t}.
\]
The radial projection is $1$-Lipschitz and commutes with orthogonal
rotations.  Moreover $\norm T_F^2\le H_R(U)$,
$\max_j\norm{V_{j\cdot}}_2\le t$, and $\norm V_F\le\sqrt K$.
Define
\[
 \beta_i=W_i\left(\sqrt{D_i}+Cd/\sqrt n\right).
\]
By \eqref{restore:tail-floor-operator}, $\beta_i^2\le Cd$.
Equation \eqref{restore:tail-weight-change} yields
\[
 \norm{\Gamma_iV}_F
 \le Ct\frac{\sqrt{D_i}}d+C\frac{d+D_i}{nd},
\]
and therefore
\begin{equation}
 \frac{C\beta_i}d\norm{\Gamma_iV}_F
 \le Ct\left\{\frac1d+\frac{\sqrt{D_i}}{d^{3/2}}\right\}.
 \label{restore:tail-bounded-weight}
\end{equation}
For the global term, use $t\ge n^{-1/2}$, $D_i\le n$, and
$d\le Cn$: $D_i/(nd^{3/2})\le t\sqrt{D_i}/d^{3/2}$ and
$1/(n\sqrt d)\le Ct/d$.
Unlike the old linear cap, the square-root cap does not give
$W_iD_i\le Cd$.  We do not impose such a bound.  Instead, the
right side of \eqref{restore:tail-bounded-weight} can exceed a
sufficiently small fixed multiple of $t$, for large $d$, only if
$D_i\ge c_*d^3$.

The unbounded part satisfies the summable estimate
\begin{equation}
 \sum_i\beta_i^2\norm{\Gamma_iT}_F^2\le C\norm T_F^2.
 \label{restore:tail-summable-weight}
\end{equation}
Indeed, after interchanging $i,j$, the local coefficient of
$\norm{T_{j\cdot}}_2^2$ is at most
$CdD_j/G_j^2\le C$.
The global coefficient is bounded by
\[
 \frac{C}{n^2d}\sum_i(d+D_i)^2\le C,
\]
because $\sum_iD_i\le Cnd$ and
$\sum_iD_i^2\le n\sum_iD_i$.
These estimates do not use a maximum-degree event.

\medskip\noindent\textbf{Independent innovations and the recurrence.}
Define the unrotated, training-measurable-coefficient innovation
\begin{equation}
 K_i=(A_{iJ}-\Omega_{iJ})W^{(-i)}f_t(U_J^{(-i)}).
 \label{restore:tail-innovation}
\end{equation}
Let $E_0=W(A-\Omega)W$.
The eigenvector equation is
\[
 U=W\Omega WU\Xi^{-1}+E_0V\Xi^{-1}+E_0T\Xi^{-1}.
\]
The first term has every row bounded by $C/\sqrt n$, by
\eqref{restore:tail-population-rows}.
For the second, compare $W_JV_J$ to
$W^{(-i)}f_t(U_J^{(-i)})O_i$ and use radial Lipschitzness,
\eqref{restore:tail-loo-residual}, and
\eqref{restore:tail-weight-change}.  This gives
\begin{align}
 \norm{U_{i\cdot}}_2\le{}&
 C/\sqrt n+(C/d)\norm{K_i}_2
 +(C\beta_i/d)\norm{U_{i\cdot}}_2\nonumber\\
 &+(C\beta_i/d)\norm{\Gamma_iV}_F
 +(C\beta_i/d)\norm{\Gamma_iT}_F
 +(C/d)\norm{(E_0T)_{i\cdot}}_2.
 \label{restore:tail-row-recursion}
\end{align}
The completed-mean diagonal contributes at most $Cpt/d\le Ct/n$
and is absorbed in the bounded terms; it is not treated as an
observed self-loop.
Although $O_i$ may depend on the tested row, the argument only uses
$\norm{K_iO_i}_2=\norm{K_i}_2$.
No conditional concentration is applied to rotated coordinates.

Take $R$ sufficiently large and then $d$ sufficiently large to absorb
the baseline and self terms.  Retain from
\eqref{restore:tail-bounded-weight} the excess
$Ct\sqrt{D_i/d^3}\ind{D_i\ge c_*d^3}$.
For $\norm{U_{i\cdot}}_2>t$, separate the cases
$\norm{K_i}_2>c'dt$ and $\norm{K_i}_2\le c'dt$ and square.
The resulting bound is
\begin{align}
 \norm{U_{i\cdot}}_2^2\ind{\norm{U_{i\cdot}}_2>t}
 \le{}&\frac C{d^2}\norm{K_i}_2^2
                         \ind{\norm{K_i}_2>c'dt}\nonumber\\
 &+\frac C{d^2}\norm{(E_0T)_{i\cdot}}_2^2
 +\frac{C\beta_i^2}{d^2}\norm{\Gamma_iT}_F^2
 +\frac{Ct^2}{d^3}D_i\ind{D_i\ge c_*d^3}.
 \label{restore:tail-pointwise-recursion}
\end{align}
Summing, using $\norm{E_0}_{\op}^2\le Cd$ and
\eqref{restore:tail-summable-weight}, and multiplying by
$\ind{\mathcal E_r}$ gives
\begin{align}
 H_R(U)\ind{\mathcal E_r}
 \le{}&\frac C{d^2}\sum_i\norm{K_i}_2^2
                         \ind{\norm{K_i}_2>c'dt}
 +\frac CdH_R(U)\ind{\mathcal E_r}\nonumber\\
 &+\frac{Ct^2}{d^3}\sum_iD_i\ind{D_i\ge c_*d^3}.
 \label{restore:tail-energy-recursion}
\end{align}
Both energy terms use the same threshold $R$.

\medskip\noindent\textbf{Expectation bounds.}
Conditional on $A_{JJ}$, \eqref{restore:tail-innovation} is a sum
of independent centered vectors.  In each coordinate its summands
are bounded by $t$, and its total variance is at most
\[
 Cp\norm{W^{(-i)}f_t(U^{(-i)})}_F^2\le Cd/n.
\]
These bounds hold for every training graph, without conditioning on
$\mathcal E_r$.  Bennett's inequality in the fixed $K$ coordinates,
followed by integration of the squared-norm tail, gives
\begin{equation}
 \E\!\left[\norm{K_i}_2^2
                        \ind{\norm{K_i}_2>c'dt}\right]
 \le Cd^2t^2e^{-c_1d\log R}.
 \label{restore:tail-innovation-moment}
\end{equation}
To see the exponent, set $v=Cd/n$ and
$h(u)=(1+u)\log(1+u)-u$.
At $x=c'dt$, the Bennett exponent
$(v/t^2)h(tx/v)$ is at least $c_1d\log R$, since
$tx/v$ is proportional to $R^2$.
Beyond that point the derivative of the exponent is at least
$c\log R/t$, which gives the same bound upon tail integration.

Every $D_i$ is a Bernoulli sum with mean at most $Cd$.
Its Chernoff bound and tail integration yield
\begin{equation}
 \E[D_i\ind{D_i\ge c_*d^3}]
 \le Cd^3e^{-c_2d^3\log d}.
 \label{restore:tail-large-degree-moment}
\end{equation}
The left side is zero if $c_*d^3>n$.
Taking expectations in \eqref{restore:tail-energy-recursion},
absorbing $C/d$, and using $nt^2=R^2$ now gives
\[
 \E[H_R(U)\ind{\mathcal E_r}]
 \le CR^2\{e^{-c_1d\log R}+e^{-c_2d^3\log d}\}.
\]
For each fixed $R$, the second exponential is eventually bounded by
the first after reducing $c_1$.  This proves
\eqref{restore:tail-energy}.  The proof is uniform under deletion
of a fixed number of vertices and uses their recomputed weights,
so it also proves the asserted training-graph version.
More generally, the same energy bound holds on any event supplying
the stated deterministic aggregate geometry, with constants depending
on those geometry bounds. The innovation and degree-tail estimates
are unconditional before restriction to that event. This observation
permits the exponentially accurate event of
\cref{restore:high-aggregate} in the higher-density part of the final
risk proof; it does not reclassify a polynomial failure probability
as an exponential one.
\end{proof}

\begin{corollary}[Exceptional coefficients are hit with exponentially small mass]
\label{restore:tail-hit-corollary}
On a training graph with the preceding floor, put
$H=G^{-1/2}UU^\top G^{1/2}$.
Let $f$ be any training-measurable vector satisfying
$\norm f_\infty\le F_0$, and let $e=(H-I)f$.
For every prescribed fixed $H_0>0$, there is a fixed $M<\infty$
such that, for any independent tested Bernoulli row with probabilities
$p_j\le Cd/n$,
\begin{equation}
 \E\!\left[\ind{\mathcal E_r}
                 \sum_{j:|e_j|>M}p_j\right]
 \le C_{H_0}e^{-H_0d}.
 \label{restore:tail-hit-mass}
\end{equation}
Here $\mathcal E_r$ is the training-measurable geometry event.
No exchangeability among degree parameters is required.
\end{corollary}

\begin{proof}
On that event, $G_j\ge a_-d$ and $\sum_jG_j\le C_Gnd$.
Consequently
\[
 |e_j|
 \le F_0+\frac{\norm{U_{j\cdot}}_2}{\sqrt{G_j}}
                  \norm{U^\top G^{1/2}f}_2
 \le F_0+F_0\sqrt{C_G/a_-}\sqrt n\norm{U_{j\cdot}}_2.
\]
For $M_R=F_0+F_0\sqrt{C_G/a_-}R$, the event
$|e_j|>M_R$ implies $\sqrt n\norm{U_{j\cdot}}_2>R$.
Thus \cref{restore:tail-energy-lemma} gives
\begin{equation}
 \E\!\left[\ind{\mathcal E_r}
                    \sum_j\ind{|e_j|>M_R}\right]
 \le\frac n{R^2}\E[H_R(U)\ind{\mathcal E_r}]
 \le Cne^{-cd\log R}.
 \label{restore:tail-coefficient-count}
\end{equation}
Multiplication by $Cd/n$ bounds the expected hit mass by
$Cd e^{-cd\log R}$.
Choose a fixed $R$ with $c\log R>H_0+1$ and absorb the factor $d$.
Independence is required only to interpret the conditional probability
that the tested row hits the training exceptional set; no independence
between the training geometry event and its eigenvectors is asserted.
\end{proof}

\begin{corollary}[Average exceptional target-row mass]
\label{restore:tail-target-corollary}
For every fixed degree cutoff $L$ and every fixed $H_0>0$, there is
a fixed $L_0$ such that
\begin{equation}
 \frac1n\sum_i\PP\!\left\{
    \sqrt{G_i}\norm{U_{i\cdot}}_2>L_0\sqrt{d/n},\quad
    D_i\le Ld,\quad \mathcal E_r\right\}
 \le C_{H_0,L}e^{-H_0d}.
 \label{restore:tail-target-mass}
\end{equation}
\end{corollary}

\begin{proof}
On the displayed event, $G_i\le C_Ld$.
Choose $L_0=R\sqrt{C_L}$.
The target-row exception then implies
$\sqrt n\norm{U_{i\cdot}}_2>R$.
Apply \eqref{restore:tail-energy} and divide its corresponding
expected count by $n$.  Choosing fixed $R$ sufficiently large gives
\eqref{restore:tail-target-mass}.
\end{proof}

\begin{remark}[The sharp-transfer consequence and its precise scope]
\label{restore:tail-scope}
Combine \eqref{restore:tail-hit-mass} with
\cref{lem:sparse-truncated-transfer} and the already established
tilted bias and quadratic estimates
\eqref{eq:floor-tilt-moments}.
For the independent training-score comparison this gives
\[
 C d^{-1/2}e^{-I}+C_{H_0}e^{-H_0d}+\PP(\mathcal E_r^c).
\]
Thus the exceptional-coefficient input is now proved for the actual
floor eigenvectors, rather than assumed.
When $d=o(\log n)$, \eqref{restore:tail-geometry-probability} is
negligible relative to $d^{-1/2}e^{-C_Id}$.
Uniformly over $d\le C_0\log n$, the same is true after taking a
fixed $r>C_IC_0+1$.
These observations introduce no logarithmic lower bound.
They concern the independent-row transfer; a complete actual-decoder
theorem also uses its separate scalar full-to-LOO comparison and
estimated-parameter argument.
\end{remark}

\subsection{Local full-to-LOO stability with an empirical degree floor}
\label{restore:loo}

The row bound used in the dense-range LOO argument need only hold on
the support of the tested row. This distinction removes a union over
all vertices. The threshold is recomputed after deletion.

\begin{lemma}[A local floored scalar comparison]\label{restore:loo-scalar}
Write $A=\left(\begin{smallmatrix}0&x^\top\\x&B\end{smallmatrix}\right)$,
$k=\sum_jx_j\le Ld$, and let $\Delta$ be the diagonal degree floor
recomputed from $B$. Suppose
\[
 cd\le\Delta_j,\qquad \sum_j\Delta_j\le Cnd,
 \qquad S_B=\Delta^{-1/2}B\Delta^{-1/2}
   =U_B\Lambda_BU_B^\top+N,
\]
where $NU_B=0$, $\norm N_{\op}\le C/\sqrt d$, and the $K$ retained
eigenvalues have a fixed gap from zero. Let the full floor be
$G=\diag(g,\Delta+E)$, with
\[
 cd\le g\le Cd,\qquad E=E_\mathrm{loc}+E_\mathrm{glob},\qquad
 0\le E_\mathrm{loc}\le\diag(x),\quad
 \norm{E_\mathrm{glob}}_{\op}\le Cd/n.
\]
Here $d\to\infty$ and $d\le Cn$. Suppose only on hit coordinates that
\begin{equation}\label{restore:loo-hit}
 \norm{\Delta_j^{-1/2}(U_B)_{j\cdot}}_2
       \le R/\sqrt{nd}\quad\text{whenever }x_j=1.
\end{equation}
Let $U,\Lambda$ be the full informative eigensystem of
$G^{-1/2}AG^{-1/2}$, and put
\[
 H_G=G^{-1/2}UU^\top G^{1/2},\qquad
 H_B=\Delta^{-1/2}U_BU_B^\top\Delta^{1/2}.
\]
For every $\norm w_\infty\le W$, uniformly over the stated inputs,
\begin{equation}\label{restore:loo-main}
 (AH_Gw)_i=x^\top H_Bw_{-i}+O_{L,R,W}(1).
\end{equation}
No bound is assumed on unhit rows of $U_B$ or on $\max_j\Delta_j$.
\end{lemma}

\begin{proof}
For $k=0$ both row actions vanish. For $k>0$, put
$F=\Delta^{-1/2}E\Delta^{-1/2}$, $h=\Delta^{-1/2}x$,
$Q=I-U_BU_B^\top$. Then $\norm F\le C/d$ and
$\norm h^2\le Ck/d$. The full normalized matrix differs from
$\diag(0,S_B)$ by $O(d^{-1/2})$: its off-diagonal row has norm
at most $C\sqrt{k/(gd)}$, and the lower renormalization has norm
$O(d^{-1})$. Thus the full signal remains separated.

Write a full generalized eigenvector as
\[
 G^{-1/2}u_r=\binom{t_r}{\Delta^{-1/2}(U_By_r+z_r)},
 \qquad U_B^\top z_r=0.
\]
The exact top and complementary equations are
\begin{align*}
 h^\top(U_By_r+z_r)&=\mu_r g t_r,\\
 (\mu_rI-N+\mu_rQFQ)z_r
   &=t_rQh-\mu_rQFU_By_r.
\end{align*}
The inverse $T_r$ of the complementary operator is bounded. Therefore
\[
 (\mu_rg-(Qh)^\top T_rQh)t_r
  =(U_B^\top h)^\top y_r
       -\mu_r(Qh)^\top T_rQFU_By_r.
\]
The denominator has magnitude at least $cd-Ck/d\ge c'd$.
It is the floored degree $g$, not the raw degree $k$, in this formula.
The matrix $Y=(y_1,\ldots,y_K)$ and its inverse have bounded norm,
by the signal-space angle and $(I+F)^{1/2}=I+O(d^{-1})$.

The local row premise and the global threshold perturbation give
\[
 \norm{U_B^\top h}\le C_Rk/\sqrt{nd},\qquad
 \norm{FU_B}_F\le C_R\sqrt{k}/(d\sqrt n)+C/n.
\]
The $C/n$ term is retained; it need not be supported on the neighbors.
Substitution yields, for $Z_\perp=(z_1,\ldots,z_K)$,
\begin{equation}\label{restore:loo-refined}
 \norm t\le C_R/\sqrt{nd},\qquad
 \norm{Z_\perp}\le C_R\{\sqrt{k}/(d\sqrt n)+1/n\}.
\end{equation}
Let $V=\Delta^{-1/2}U_B$, $v_r=\Delta^{-1/2}z_r$ and
$f_r=Vy_r+v_r$. Exact generalized orthonormality gives
\[
 Y^\top Y+Z_\perp^\top Z_\perp+gtt^\top+f^\top Ef=I.
\]
On the local part of $E$, use the hit-row bound; on its global part
use $\norm{E_\mathrm{glob}}\le Cd/n$ and
$\norm f\le C/\sqrt d$. Consequently
\begin{equation}\label{restore:loo-gram}
 \norm{YY^\top-I}\le C_R\{k/(nd)+1/n\}.
\end{equation}

For $w_V=V^\top\Delta w_{-i}$, the exact row expansion is
\[
 (AH_Gw)_i=
 \sum_r(x^\top Vy_r+x^\top v_r)
 \{g t_rw_i+y_r^\top w_V+v_r^\top\Delta w_{-i}
                                      +f_r^\top Ew_{-i}\}.
\]
Only total degree is needed to bound global coefficient vectors. In
particular,
\begin{align*}
 \norm{x^\top V}&\le C_Rk/\sqrt{nd},&
 \norm{w_V}&\le C_W\sqrt{nd},\\
 \norm{x^\top v}&\le C_R\{k/(d^{3/2}\sqrt n)+\sqrt{k}/(n\sqrt d)\},&
 \norm{v^\top\Delta w_{-i}}&\le C_{R,W}\{\sqrt{k/d}+\sqrt{d/n}\},\\
 \norm{gt}&\le C_R\sqrt{d/n},&
 \norm{f^\top Ew_{-i}}&\le C_{R,W}\{k/\sqrt{nd}+\sqrt{d/n}\}.
\end{align*}
For example the global part of the last bound follows from
$|f_r^\top E_\mathrm{glob}w_{-i}|
\le (Cd/n)W\sqrt n\norm{f_r}\le C_W\sqrt{d/n}$.
The leading term is $x^\top VYY^\top w_V$; its error relative to
$x^\top VV^\top\Delta w_{-i}$ is
$O(k^2/(nd)+k/n)$. Every other product is bounded by a fixed
multiple of
\[
 k/d+\sqrt{k/n}+k^2/(nd)+k/n+k\sqrt{k}/(d\sqrt n).
\]
These quantities are bounded for $k\le Ld$ and $d\le Cn$.
This proves the scalar comparison.
\end{proof}

For the empirical floor $\tau=a\bar D$, the displayed decomposition of
$E$ is automatic on the two-sided total-degree event. Indeed, if
$\bar D_B$ is the mean degree of $B$, then
\[
 \tau_A-\tau_B=a(2k-\bar D_B)/n.
\]
First change each deleted degree by $x_j$ while keeping $\tau_B$;
this yields $E_\mathrm{loc}\in[0,\diag(x)]$. Then change the threshold;
this yields $\norm{E_\mathrm{glob}}\le C_Ld/n$.

\begin{corollary}[Exact-degree reconstruction and simple hollowing]
\label{restore:loo-corrected}
Put $R_0=AH_G$ and $\mathbf d=A\one$, and define
\begin{equation}\label{restore:reconstruction}
 R_1=R_0+(\mathbf d-R_0\one)\one^\top/n,\qquad
 R^h=R_1+\frac{\diag(R_1)}{n-1}\one^\top
                -\frac n{n-1}\diag\{\diag(R_1)\}.
\end{equation}
In this display $\diag(R_1)$ in the numerator denotes its diagonal
column vector; the final $\diag\{\cdot\}$ forms a diagonal matrix.
Then $R^h\one=\mathbf d$ and $\diag(R^h)=0$ exactly.
For $\bar w_{-i}=(n-1)^{-1}\sum_{j\ne i}w_j$, the same local premises
give
\begin{equation}\label{restore:loo-corrected-action}
 (R^h w)_i
 =x^\top\{H_B(w_{-i}-\bar w_{-i}\one)+\bar w_{-i}\one\}
       +O_{L,R,W}(1).
\end{equation}
\end{corollary}

\begin{proof}
For $\bar w=n^{-1}\sum_jw_j$, exact degree correction gives
$(R_1w)_i=(R_0(w-\bar w\one))_i+k\bar w$.
Apply the preceding lemma. On hit coordinates $|(H_B\one)_j|\le C_R$,
because $\norm{U_B^\top\Delta^{1/2}\one}\le C\sqrt{nd}$.
Replacing $\bar w$ by $\bar w_{-i}$ therefore costs $O(k/n)$.
Moreover \eqref{restore:loo-refined} gives
$\norm{u_i}^2=g\norm t^2\le C_R/n$, hence
$|(R_0)_{ii}|\le C_Rd/n$. Also $(R_0\one)_i=O_R(k+1)$.
Thus $|(R_1)_{ii}|\le C_Rd/n$ and the hollow change to a bounded
contrast is at most $2W|(R_1)_{ii}|=O(d/n)$. The two exact identities
follow by summing the defined rows and examining their diagonal.
\end{proof}

The independent coefficient vector in
\eqref{restore:loo-corrected-action} has perturbation
$(H_B-I)(w_{-i}-\bar w_{-i}\one)$. For block-constant contrasts it
therefore has the same tilted cancellation and quadratic bound already
proved in \eqref{eq:floor-tilted-cancellation}--\eqref{eq:floor-tilt-moments}.
For clarity, the hit-row exception is paid directly. On the
training event, $\Delta_j\ge c d$, so
\[
 \mathcal J_{\rm bad}
 :=\{j:\norm{\Delta_j^{-1/2}(U_B)_{j\cdot}}_2
                          >R/\sqrt{nd}\}
 \subseteq
 \{j:\sqrt n\norm{(U_B)_{j\cdot}}_2>R\sqrt c\}.
\]
\Cref{restore:tail-energy-lemma} bounds the expected cardinality of
this set on the training event by $Cn e^{-c'd\log R}$.
Conditioning only on the deleted graph and using $p_{ij}\le Cd/n$
therefore bounds the probability that the tested row hits it by
$Cd e^{-c'd\log R}$. Choosing a sufficiently large fixed $R$
gives any prescribed fixed exponential rate $H$.
The upper-degree exception $k>Ld$ has arbitrarily large exponential
rate after increasing the fixed $L$. Neither bound uses a union over
all possible tested targets.

Uniformity in $w$ is an algebraic stability statement. It does not make
full-graph fitted coefficients independent of $x$; the following fitted
parameter lemma replaces them by deterministic true contrasts before
the independent sharp-tail argument is used.

\subsection{Sparse DC fitting from aggregate geometry and a local row bound}
\label{restore:fit}

This subsection supplies the deterministic fitting module for the
degree-floor decoder. It uses the exact Bernoulli working criterion and
the exact-degree reconstruction $R^h$ in \eqref{restore:reconstruction}.
In particular, its rank-one degree correction and simple row hollowing
are different from the earlier correction $\diag(s)(H-I)$.
One profile refinement suffices. The optional finite repeated hard-refit
version is covered explicitly by \cref{restore:fit-invariant-basin}.

Put $d=np$, $B=pC$, $\Omega=\Theta ZBZ^\top\Theta$ and
$D_{0,i}=(\Omega\one)_i$. As before, $\mathbf d=A\one$ is the raw
degree vector, $D_i=\mathbf d_i$, and
\[
 G_i=\max(D_i,a_0\overline D),\qquad
 S=G^{-1/2}AG^{-1/2},\qquad
 R_0=G^{1/2}U\Lambda U^\top G^{1/2}.
\]
The constants $a_0>0$ and all subsequent clipping constants are fixed
using only the structural parameter bounds. They do not depend on $n$,
the unknown labels, or individual unobserved parameters.
Work on the two-sided total-degree event on which
$\overline D\asymp d$ and the population floor is inactive.
The aggregate geometric and block-total premises used below are
\begin{align}
 \norm{S-G^{-1/2}\Omega G^{-1/2}}_{\op}
   &\le C/\sqrt d,&
 \sum_iG_i&\le Cnd, \label{restore:fit-geometry-one}\\
 \sum_i\frac{(G_i-D_{0,i})^2}{G_i}
 +\sum_i\frac{(D_i-D_{0,i})^2}{G_i}
   &\le Cn,&
 \norm{(I-UU^\top)G^{1/2}Z}_F&\le C\sqrt n,
 \label{restore:fit-geometry-two}\\
 \max_{a,b}\left|[Z^\top(A-\Omega)Z]_{ab}\right|
   &\le Cn.&& \label{restore:fit-block-event}
\end{align}
Here $G_i\ge cd$, $D_{0,i}\asymp d$, and the retained signal has a fixed
gap from zero. The degree-floor proposition already supplies
\eqref{restore:fit-geometry-one}--\eqref{restore:fit-geometry-two}.
The fixed-$K$ Bernoulli block-total bound
\eqref{restore:fit-block-event} has failure probability
$C\exp\{-cL^2n/(d+L)\}$ when its fixed threshold constant is $L$.
Under $np^2=d^2/n\le C_1$, increasing $L$ pays this failure with any
prescribed fixed exponential rate in $d$. No degree-wise union bound
is needed for this block event.

Write $\Pi=UU^\top$. An aligned orthonormal basis of the random
population space has the exact form
\[
 V=G^{-1/2}\Theta Z
       (Z^\top\Theta G^{-1}\Theta Z)^{-1/2}O,
\]
where $O$ is orthogonal. The diagonal Gram matrix has entries comparable
to $n/d$. The signal gap and \eqref{restore:fit-geometry-one} imply
\begin{equation}\label{restore:fit-basis}
 \norm{U-V}_F\le C/\sqrt d,\qquad
 \sqrt{G_i}\norm{V_{i\cdot}}_2\le C\sqrt{d/n}.
\end{equation}
An orthogonal rotation used in this display does not change any row
norm or projection. Also $\norm S_{\op}\le1$: the inequality
$|x^\top Ax|\le\sum_iD_ix_i^2\le\sum_iG_ix_i^2$ proves this directly.
Hence $\norm\Lambda_{\op}\le1$.

\begin{lemma}[Uniform weighted mass of wrong-label sets]
\label{restore:fit-wrong-mass}
Fix $M<\infty$. For every possibly data-dependent set $E$ with
$|E|\le Mn/d$, on the aggregate event,
\begin{equation}\label{restore:fit-wrong-mass-bounds}
 \sum_{i\in E}G_i\le C_Mn,\qquad
 \sum_{i\in E}\sqrt{G_i}\norm{U_{i\cdot}}_2
       \le C_M\sqrt{n/d}.
\end{equation}
Moreover
$\sum_i\sqrt{G_i}\norm{U_{i\cdot}}_2\le C\sqrt{nd}$.
\end{lemma}

\begin{proof}
The elementary inequality
$G_i\le2D_{0,i}+(G_i-D_{0,i})^2/G_i$ gives
$\sum_{i\in E}G_i\le C(d|E|+n)\le C_Mn$.
Put $a_i=\sqrt{G_i}\norm{U_{i\cdot}}_2$.
By \eqref{restore:fit-basis} and Cauchy--Schwarz,
\[
 \sum_{i\in E}a_i
 \le C|E|\sqrt{d/n}
      +\left(\sum_{i\in E}G_i\right)^{1/2}\norm{U-V}_F
 \le C_M\sqrt{n/d}.
\]
The last assertion follows from
$\sum_i a_i\le(\sum_iG_i)^{1/2}\norm U_F$ and fixed $K$.
All estimates hold simultaneously for all such sets; they impose no
independence requirement on the selected labels.
\end{proof}

\begin{lemma}[Radius-count initialization under aggregate geometry]
\label{restore:fit-initializer}
The deterministic radius-count initializer specified earlier has at most
$C n/d$ errors on the aggregate event, after one common permutation.
This assertion requires neither a uniform degree upper bound nor
uniform spectral-row incoherence.
\end{lemma}

\begin{proof}
The normalized rows of $V$ are the $K$ orthogonal population centers.
Choose fixed $b_0>2\max_iD_{0,i}/d$, uniformly over the structural class.
For every $G_i>b_0d$,
$(G_i-D_{0,i})^2/G_i\ge G_i/4\ge b_0d/4$.
Thus at most $Cn/d$ such rows occur by
\eqref{restore:fit-geometry-two}. They contribute at most $4Cn/d$
to the squared directional error. For the remaining rows,
$\norm{V_{i\cdot}}_2\ge c/\sqrt n$. The row-normalization inequality
and \eqref{restore:fit-basis} bound their total squared directional
error by $Cn\norm{U-V}_F^2\le Cn/d$.
The existing radius-count lemma now applies with total squared error
$Cn/d$ and fixed positive community proportions. Its proof is
deterministic and unchanged.
\end{proof}

\begin{lemma}[Degree-corrected block totals and absolute mass]
\label{restore:fit-block-totals}
For $R^h$ as in \eqref{restore:reconstruction},
\begin{equation}\label{restore:fit-global-absolute}
 \sum_{i,j}|R^h_{ij}|\le Cnd.
\end{equation}
Uniformly over every fitted partition $\widehat Z$ having at most
$Mn/d$ errors under the common label matching,
\begin{equation}\label{restore:fit-total-accuracy}
 [\widehat Z^\top R^h\widehat Z]_{ab}
       =n_an_bB_{ab}+O_M(n).
\end{equation}
The omitted self-loops are included in the remainder.
\end{lemma}

\begin{proof}
Let $S_{\rm tail}=(I-\Pi)S(I-\Pi)$ and
$F_1=(I-\Pi)G^{1/2}\one$. Spectral invariance and the signal gap give
$\norm{S_{\rm tail}}_{\op}\le C/\sqrt d$ and
$\norm{F_1}_2\le C\sqrt n$.
Writing $v=\mathbf d-R_0\one$, one has exactly
$v=G^{1/2}S_{\rm tail}F_1$. Hence
\[
 \sum_i|v_i|
 \le\left(\sum_iG_i\right)^{1/2}
          \norm{S_{\rm tail}F_1}_2\le Cn.
\]
Since $\Pi_{ii}\le1$ and $\sum_i\Pi_{ii}=K$,
\[
 \sum_i|(R_0)_{ii}|
 \le\sum_iG_i\Pi_{ii}
 \le2(\max_iD_{0,i})K
       +\sum_i\frac{(G_i-D_{0,i})^2}{G_i}
 \le C(d+n)\le Cn.
\]
Consequently $\sum_i|(R_1)_{ii}|\le Cn$ and the simple hollowing
has total absolute cost at most $2\sum_i|(R_1)_{ii}|\le Cn$.
With $a_i$ as in the preceding lemma,
$|(R_0)_{ij}|\le a_ia_j$, so
$\sum_{i,j}|(R_0)_{ij}|\le(\sum_i a_i)^2\le Cnd$.
The rank-one correction has total absolute mass $\sum_i|v_i|\le Cn$.
This proves \eqref{restore:fit-global-absolute}.

Put $\widehat F=(I-\Pi)G^{1/2}\widehat Z$. By
\eqref{restore:fit-wrong-mass-bounds},
\[
 \norm{\widehat F}_F
 \le\norm{(I-\Pi)G^{1/2}Z}_F
       +\norm{G^{1/2}(\widehat Z-Z)}_F
 \le C_M\sqrt n .
\]
The exact quadratic identity
\[
 \widehat Z^\top(R_0-A)\widehat Z
       =-\widehat F^\top S_{\rm tail}\widehat F
\]
therefore has norm at most $C_Mn/\sqrt d$. Moving the mislabeled
vertices changes a raw block total by at most a fixed multiple of
$\sum_{i\in E}D_i\le\sum_{i\in E}G_i\le C_Mn$.
Equation \eqref{restore:fit-block-event} and the exact within-block
mean normalization of $\theta$ prove \eqref{restore:fit-total-accuracy}
with $R_0$ in place of $R^h$.
The rank-one correction changes a fitted block total by
$[\widehat Z^\top v]_a\widehat n_b/n$, where
$\widehat Z^\top v=\widehat F^\top S_{\rm tail}F_1=O_M(n/\sqrt d)$.
The already bounded hollowing cost completes the proof.
\end{proof}

\subsubsection{A bounded, normalized degree-parameter estimator}

Choose fixed positive $\alpha$ sufficiently small and $\beta$ sufficiently
large, using the structural bounds, so that on the two-sided
total-degree event
\[
 \alpha\overline D\le\tfrac12\min_iD_{0,i},\qquad
 \beta\overline D\ge2\max_iD_{0,i}.
\]
Define
\begin{equation}\label{restore:fit-clamped-theta}
 T_i=\min\{\max(D_i,\alpha\overline D),\beta\overline D\},
 \qquad
 \widehat\theta_i=
 \frac{T_i}{\widehat n_{\widehat z_i}^{-1}
                 \sum_{j:\widehat z_j=\widehat z_i}T_j}.
\end{equation}
This changes the nuisance estimator only. The retained degrees and the
degree identity $R^h\one=\mathbf d$ still use the original raw degrees.
The fitted parameters lie in $[\alpha/\beta,\beta/\alpha]$, and their
mean in every fitted block is exactly one.

\begin{lemma}[Clamped nuisance estimates]
\label{restore:fit-theta}
Uniformly over the $Mn/d$ count basin,
\begin{equation}\label{restore:fit-theta-rates}
 \frac{\norm{\widehat\theta-\theta}_2}{\sqrt n}
        \le C_Md^{-1/2},\qquad
 \widehat n_a^{-1}\sum_{\widehat z_i=a}T_i
        =dS_a+O_M(1),
\end{equation}
where $D_{0,i}=dS_{z_i}\theta_i$.
\end{lemma}

\begin{proof}
The chosen clamp interval contains $D_{0,i}$ with a fixed margin.
A direct split into the unclamped, lower-clamped, and upper-clamped
cases gives
\[
 |T_i-D_i|\le C\frac{(D_i-D_{0,i})^2}{G_i},\qquad
 |T_i-D_{0,i}|^2\le Cd\frac{(D_i-D_{0,i})^2}{G_i}.
\]
For the first inequality, an active lower clamp implies a degree
deviation of order $d$, and an active upper clamp implies a deviation
of order $D_i$. For the second, the unclamped range has $G_i\le Cd$;
projection onto an interval containing $D_{0,i}$ cannot increase
the deviation. Above the upper clamp,
$D_i((T_i-D_{0,i})/(D_i-D_{0,i}))^2\le T_i\le Cd$.
Summing yields
\begin{equation}\label{restore:fit-clamp-cost}
 \sum_i|T_i-D_i|\le Cn,\qquad
 \sum_i|T_i-D_{0,i}|^2\le Cnd.
\end{equation}
Since $\one=Z\one_K$, the raw block-total event implies
$\sum_{z_i=a}(D_i-D_{0,i})=O(n)$.
Clamping changes that sum by $O(n)$, and moving $Mn/d$ labels changes
a clamped degree total by at most $Cd(Mn/d)=O_M(n)$.
The fitted size differs from $n_a$ by $O_M(n/d)$.
Together with $\sum_{z_i=a}D_{0,i}=n_adS_a$, these facts prove the
group-mean assertion in \eqref{restore:fit-theta-rates}.
For correctly labeled vertices, division of
$T_i-D_{0,i}$ by this fitted mean, followed by
\eqref{restore:fit-clamp-cost}, gives squared error $O_M(n/d)$.
The incorrectly labeled vertices number $O_M(n/d)$ and have bounded
theta error. This proves the first assertion.
\end{proof}

\subsubsection{The exact Bernoulli block likelihood}

Fit a positive symmetric $B$ by maximizing the ordered Bernoulli
working criterion with $R^h$ and \eqref{restore:fit-clamped-theta},
divided by two, under the same fixed probability cap
$\widehat\theta_i\widehat\theta_jB_{\widehat z_i\widehat z_j}
\le q_{\max}<1$.
The next argument uses aggregate absolute mass instead of the stronger
and unnecessary assertion that every reconstructed entry is less than one.

\begin{lemma}[Unique Bernoulli block fit and its rate]
\label{restore:fit-block-root}
The fitted block roots are unique and interior for all large $n$.
Uniformly over the $Mn/d$ basin,
\begin{equation}\label{restore:fit-B-rate}
 \max_{a,b}\left|\frac{\widehat B_{ab}}{B_{ab}}-1\right|
       \le C_M(d^{-1}+p)\le C_Md^{-1}.
\end{equation}
\end{lemma}

\begin{proof}
For the ordered pairs $\mathcal J_{ab}$ associated with an unordered
block pair, put
\[
 E_{ab}=\sum_{\mathcal J_{ab}}R^h_{ij},\qquad
 t_{ij}=\widehat\theta_i\widehat\theta_j,\qquad
 T_{ab}=\sum_{\mathcal J_{ab}}t_{ij},\qquad
 \nu_{ab}=2-\mathbf1\{a=b\}.
\]
The preceding block-total lemma and exact fitted normalization imply
\[
 E_{ab}=\nu_{ab}n_an_bB_{ab}+O_M(n)\asymp n^2p>0,\qquad
 T_{ab}=\nu_{ab}\widehat n_a\widehat n_b+O(n).
\]
The exact score equation is $E_{ab}=F_{ab}(b)$, where
\[
 F_{ab}(b)=b\sum_{\mathcal J_{ab}}
              \frac{(1-R^h_{ij})t_{ij}}{1-bt_{ij}}.
\]
For every admissible $b$, bounded positive theta, the probability cap,
and \eqref{restore:fit-global-absolute} give
\[
 F_{ab}'(b)
 =\sum_{\mathcal J_{ab}}
          \frac{(1-R^h_{ij})t_{ij}}{(1-bt_{ij})^2}
 \ge cn^2-C\sum_{\mathcal J_{ab}}|R^h_{ij}|
 \ge cn^2-Cnd>0.
\]
Here $p=d/n\to0$. At zero, $F_{ab}(0)=0$; its value at the
positive cap endpoint is of order $n^2$, while $E_{ab}\asymp n^2p$.
There is a unique interior root. The likelihood derivative is
$(E_{ab}-F_{ab}(b))/b$, so the root is its unique global maximum.

For $b\asymp p$,
\[
 \sum_{\mathcal J_{ab}}
       \frac{(1-R^h_{ij})t_{ij}}{1-bt_{ij}}
       =T_{ab}+O(n^2p).
\]
The error follows by expanding the denominator and using
$\sum|R^h_{ij}|\le Cnd=Cn^2p$.
Bracketing the root between fixed multiples of $p$, then using this
expansion, gives
$\widehat B_{ab}=(E_{ab}/T_{ab})\{1+O(p)\}$.
The block-total and group-size errors prove the first inequality in
\eqref{restore:fit-B-rate}. The second uses $dp=np^2\le C_1$.
\end{proof}

\subsubsection{A local row condition suffices for every fitted stage}

For a tested vertex $i$, impose only the local conditions
\begin{equation}\label{restore:fit-local-row}
 D_i\le bd,\qquad
 a_i:=\sqrt{G_i}\norm{U_{i\cdot}}_2\le L\sqrt{d/n},
\end{equation}
where $b,L$ are fixed. These are not imposed on every neighbor
simultaneously.

\begin{lemma}[Local absolute mass and wrong-label stability]
\label{restore:fit-local-mass}
On the aggregate event and \eqref{restore:fit-local-row},
\begin{equation}\label{restore:fit-local-mass-bounds}
 \sum_j|R^h_{ij}|\le C_{b,L}d,\qquad
 \sup_{|E|\le Mn/d}\sum_{j\in E}|R^h_{ij}|\le C_{M,b,L}.
\end{equation}
\end{lemma}

\begin{proof}
Since $|(R_0)_{ij}|\le a_ia_j$,
\[
 \sum_j|(R_0)_{ij}|\le a_i\sum_j a_j\le Cd,\qquad
 \sum_{j\in E}|(R_0)_{ij}|
       \le a_i\sum_{j\in E}a_j\le C_M
\]
by \cref{restore:fit-wrong-mass}. Also
$|v_i|\le D_i+\sum_j|(R_0)_{ij}|\le Cd$.
The rank-one correction therefore adds at most $Cd$ to the first
bound and $|v_i||E|/n\le C_M$ to the second.
Finally $|(R_1)_{ii}|\le a_i^2+|v_i|/n\le Cp$.
The simple row hollowing changes either row bound by at most
$2|(R_1)_{ii}|\le Cp$. This proves the claim.
\end{proof}

\begin{lemma}[Bernoulli profile reduction for a signed row]
\label{restore:fit-profile}
Let $r$ be a hollow signed row satisfying
$\sum_jr_j=D_i\ge0$ and $\sum_j|r_j|\le Ld$.
Let the positive neighbor intensities $b_j$ lie between fixed
multiples of $p$. Maximizing the exact Bernoulli row criterion
\[
 \sum_j\{r_j\log(tb_j)+(1-r_j)\log(1-tb_j)\}
\]
over $t\ge0$ under the fixed cap $tb_j\le q_{\max}<1$ gives the
usual normalized block-profile contrast, with remainder $O(np^2)$.
No entrywise bound on $r_j$ is required.
\end{lemma}

\begin{proof}
After combining the coefficient of $\log t$, the derivative is
\[
 D_i/t-\sum_j\frac{(1-r_j)b_j}{1-tb_j}.
\]
Throughout the admissible interval,
\[
 \sum_j\frac{(1-r_j)b_j}{1-tb_j}
 \ge cnp-Cp\sum_j|r_j|
 \ge cd-CLdp\ge c'd.
\]
Thus every maximizing $t$ lies in a fixed bounded interval,
because $D_i\le Ld$. This reasoning does not require individual
nonedge weights $1-r_j$ to be nonnegative.
On that fixed interval, the exact Bernoulli expansion is uniform:
\[
 \ell_r(t)
   =D_i\log t+\sum_jr_j\log b_j
         -t\sum_jb_j+O(np^2).
\]
Restoring the edge terms in the linear exposure costs
$Cp\sum_j|r_j|=O(np^2)$, and the quadratic remainder costs
$Cp^2(n+\sum_j|r_j|)=O(np^2)$.
The leading optimizer $D_i/\sum_jb_j$ is bounded and admissible
for all large $n$, so taking maxima preserves the remainder.
The resulting label-dependent maximum is the standard normalized
block-profile score. If $D_i=0$, the derivative bound makes $t=0$
optimal after combining the $\log t$ terms; use the continuous
convention $0\log0=0$.
\end{proof}

\begin{proposition}[Uniform replacement by the true DC contrast]
\label{restore:fit-uniform-score}
\label{restore:fit-score}
Fix $M,b,L$. Use the clamped theta estimates
\eqref{restore:fit-clamped-theta}, the exact Bernoulli block fit, and
the exact Bernoulli target profile on $R^h$.
On the aggregate event and the local target condition
\eqref{restore:fit-local-row}, uniformly over every fitted partition
in the $Mn/d$ basin,
\begin{equation}\label{restore:fit-true-score}
 \widehat L_i^{ab}
       =[R^h w_{i,ab}]_i+O_{M,b,L}(1),\qquad
 (w_{i,ab})_j=
     \log\frac{q_{a,z_j}^{(i)}}{q_{b,z_j}^{(i)}}\quad(j\ne i).
\end{equation}
The true profiles $q^{(i)}$ are the self-excluded neighbor profiles
defined earlier. The value of $w_{i,ab}$ at $i$ is immaterial since
$R^h_{ii}=0$.
\end{proposition}

\begin{proof}
By \cref{restore:fit-profile}, the profile reduction costs
$O(np^2)=O(1)$. Exact fitted-block normalization gives the fitted
neighbor mass
$\widehat M_c^{(i)}
 =\widehat n_c-\widehat\theta_i\mathbf1\{\widehat z_i=c\}$.
The fitted sizes differ from the true sizes by $O_M(n/d)$,
the fitted theta values are bounded, and
\eqref{restore:fit-B-rate} gives relative block error $O_M(1/d)$.
Smoothness of normalization and logarithms on the fixed positive
profile class therefore yields fitted log-profile error $O_M(1/d)$.
Multiplication by the row absolute mass in
\eqref{restore:fit-local-mass-bounds} costs $O_M(1)$.
Replacing the fitted labels by the true labels changes a bounded
coefficient vector on at most $Mn/d$ coordinates; the second bound
in \eqref{restore:fit-local-mass-bounds} gives another $O_M(1)$.
All inequalities hold simultaneously for every partition in the basin,
so data-dependent fitting does not introduce an independence assumption.
\end{proof}

\begin{lemma}[A uniform basin for finite hard refits]
\label{restore:fit-invariant-basin}
\label{restore:fit-invariant}
There is a fixed $M$ such that, on the aggregate event and for all
sufficiently large $n$, each hard nuisance-fit and profile-assignment
update maps the $Mn/d$ error basin into itself under the same label
matching. Outside a fixed target-exception set determined by the graph,
every error at every updated iterate is contained in
\[
 [R^h w_{i,ab}]_i\le C_M
 \quad\text{for some }b\ne z_i.
\]
The target-exception set can be taken to be the failure of the local
conditions \eqref{restore:fit-local-row}, with fixed sufficiently large
thresholds. Thus a finite number of these specified hard refits, including
adaptive finite stopping after at least one update, requires no union
over iterations in the final risk bound.
\end{lemma}

\begin{proof}
Fix sufficiently large constants $b,R$. The set
$\{i:D_i>bd\}$ has at most $Cn/d$ elements by the degree energy in
\eqref{restore:fit-geometry-two}, since $b$ exceeds twice the largest
population degree multiplier. Since $G_i\ge cd$,
\eqref{restore:fit-basis} gives $\norm{V_{i\cdot}}_2\le C/\sqrt n$.
For $R>2C$, the set $\{i:\sqrt n\norm{U_{i\cdot}}_2>R\}$
also has at most $Cn/d$ elements, by
$\norm{U-V}_F^2\le C/d$.
Every other target satisfies \eqref{restore:fit-local-row} for a
fixed $L$, because $G_i\le C_bd$ there.
This target-exception set depends only on the graph and its
cardinality bound is independent of $M$.

Write $M_G=G^{-1/2}\Omega G^{-1/2}$ and
$S_K=U\Lambda U^\top$. Best rank-$K$ approximation and
\eqref{restore:fit-geometry-one} imply
$\norm{S_K-M_G}_{\op}\le C/\sqrt d$. Hence
\[
 \norm{G^{-1/2}(R_0-\Omega)Z}_F^2
 \le \frac Cd\norm{G^{1/2}Z}_F^2\le Cn.
\]
Moreover
$\norm{G^{-1/2}v}_2^2
 =\norm{S_{\rm tail}F_1}_2^2\le Cn/d$.
The rank-one correction therefore preserves the bound
\begin{equation}\label{restore:fit-basin-action}
 \norm{G^{-1/2}(R_1-\Omega)Z}_F^2\le Cn.
\end{equation}
The simple hollowing changes any bounded block-constant contrast at
row $i$ by at most $C|(R_1)_{ii}|$.
Since $\sum_i|(R_1)_{ii}|\le Cn$, at most $Cn/d$ rows can incur
a hollowing contrast error larger than any prescribed fixed small
multiple of $d$.
This auxiliary counting set is needed only for the deterministic
basin proof; it is not added as a probability-cost event in the sharp
risk argument.

The true profile contrast has a uniform positive population margin
$[\Omega w_{i,ab}]_i\ge c_{\rm gap}d$ for $a=z_i$, $b\ne a$,
including its $O(p)$ self-exclusion correction.
Extend each $w_{i,ab}$ to a block-constant vector at coordinate $i$;
this does not affect $R^h w$ because $R^h_{ii}=0$.
Its $K$ block coefficients are uniformly bounded. For a regular
target, \cref{restore:fit-uniform-score} bounds the fitted discrepancy
by $C_M$, which is less than $c_{\rm gap}d/4$ for all sufficiently
large $n$ (the threshold may depend on $M$).
Outside the auxiliary hollowing set, a new erroneous label therefore
forces
\[
 \norm{[(R_1-\Omega)Z]_{i\cdot}}_2\ge c d.
\]
For the regular targets, $G_i\le C_bd$. Summing the squared
requirement divided by $G_i$ and applying
\eqref{restore:fit-basin-action} gives at most $C_0n/d$ such targets.
Including the two cardinality exceptions still gives $C_0n/d$
with $C_0$ independent of $M$.

Choose $M$ at least both this $C_0$ and the initializer constant.
The basin is invariant for all sufficiently large $n$.
For every regular target, the canonical error-event inclusion
follows directly from \eqref{restore:fit-true-score}, independently
of this counting argument. It is the same event for all partitions
in the fixed basin. Induction therefore covers each finite hard
refit and any data-dependent finite stopping rule without a new
probability union or a stopping-time independence claim.
\end{proof}

\subsubsection{Ordinary SBM with a block Bernoulli fit and no target profiling}

The sparse ordinary-SBM decoder in this subsection sets $\theta_i=1$,
uses the same radius-count initializer and the same $R^h$, fits a
symmetric block Bernoulli matrix, and compares the full Bernoulli row
likelihoods. The target parameter is not profiled. This is an explicitly
specified block-fitting decoder; it is not a claim about the earlier
unmodified free-profile mixture E/M schedule.

\begin{proposition}[Ordinary block fitting and the true Bernoulli LRT]
\label{restore:fit-ordinary}
Assume ordinary SBM with a fixed positive well-conditioned block profile,
comparable group sizes, $d=np\to\infty$ and $np^2\le C_1$.
On the aggregate event, the radius-count initialization has $O(n/d)$
errors. For every fitted partition in a fixed $Mn/d$ basin, the symmetric
block Bernoulli working MLE based on $R^h$ satisfies
\begin{equation}\label{restore:fit-ordinary-parameter}
 \max_{a,b}|\widehat P_{ab}-P_{ab}|\le C_M/n.
\end{equation}
At every target obeying \eqref{restore:fit-local-row}, its full
Bernoulli contrast satisfies
\begin{equation}\label{restore:fit-ordinary-score}
 \widehat L_i^{ab}
    =[R^h w_{ab}]_i+c_i^{ab}+O_M(1),
\end{equation}
where
\[
 (w_{ab})_j=
   \log\frac{P_{a,z_j}(1-P_{b,z_j})}
                  {P_{b,z_j}(1-P_{a,z_j})},\qquad
 c_i^{ab}=\sum_{j\ne i}
       \log\frac{1-P_{a,z_j}}{1-P_{b,z_j}}.
\]
The counterpart of \cref{restore:fit-invariant-basin} also holds for
these ordinary block fits and full Bernoulli assignments.
\end{proposition}

\begin{proof}
For each unordered fitted block pair, let $E_{ab}$ be the corresponding
ordered reconstructed total and $N_{ab}$ its number of off-diagonal
ordered pairs. The exact Bernoulli working criterion is
$E_{ab}\log b+(N_{ab}-E_{ab})\log(1-b)$.
The block-total lemma gives
$E_{ab}=\nu_{ab}n_an_bP_{ab}+O_M(n)$ and
$N_{ab}=\nu_{ab}\widehat n_a\widehat n_b+O(n)$.
For all sufficiently large $n$, $0<E_{ab}<N_{ab}$; its unique MLE is
$\widehat P_{ab}=E_{ab}/N_{ab}$.
Any fixed sufficiently small probability floor of order $\overline D/n$
is inactive on this event. Since
$|\widehat n_a-n_a|\le Mn/d$ and $P_{ab}\asymp p$,
\[
 |\widehat P_{ab}-P_{ab}|
 \le C_M\{n/n^2+p(n/d)/n\}=C_M/n.
\]
This proves \eqref{restore:fit-ordinary-parameter}.

Writing the Bernoulli contrast as an incident-edge log-odds action
plus its nonedge intercept is exact. On the fixed positive profile
class, \eqref{restore:fit-ordinary-parameter} changes each log-odds
coefficient by $O_M((np)^{-1})=O_M(d^{-1})$.
Its total contribution at a regular target is $O_M(1)$ by the row
absolute-mass bound. Incorrect fitted neighbor labels change a bounded
log-odds coefficient on at most $Mn/d$ coordinates, at a total row
cost $O_M(1)$ by \cref{restore:fit-local-mass}.

For the intercept, the derivatives of $\log(1-b)$ are bounded on
the relevant probability range. Parameter replacement therefore costs
$n\,O_M(1/n)=O_M(1)$. Changing $Mn/d$ group assignments changes
each nonedge log-ratio by at most $Cp$, so its additional cost is
$C_M(n/d)p=O_M(1)$. These observations prove
\eqref{restore:fit-ordinary-score} for all data-dependent basin
partitions at once.

For the basin argument, the true population contrast equals a sum
of Bernoulli Kullback--Leibler divergences, up to the $O(p)$ completed
diagonal, and is at least $c_{\rm gap}d$.
The block-constant log-odds vectors are uniformly bounded.
Consequently the weighted action estimate
\eqref{restore:fit-basin-action}, the same target-exception count,
and the same hollowing count prove an output error bound $C_0n/d$.
As before $C_0$ is independent of the chosen input basin radius $M$;
only the sufficiently-large-$n$ threshold depends on $M$.
Choosing $M$ large enough proves invariance. Every regular-target
error at every refit is contained in the same true-contrast event
\eqref{restore:fit-ordinary-score} with its fixed $O_M(1)$ threshold,
so finite adaptive hard refits need no iteration union.
\end{proof}

The remaining interface is local. The radius-count initializer and
nuisance fit above need only aggregate geometry. To turn
\eqref{restore:fit-true-score} into the actual sharp risk theorem one
pays the target-local spectral-row exception, then applies the local
full-to-LOO comparison in \cref{restore:loo-corrected} to the deterministic
true contrast. The fitted coefficients themselves are not conditioned
on as if independent of the tested row.
The high-degree part of \eqref{restore:fit-local-row} has arbitrarily
large fixed exponential rate after increasing $b$. On $D_i\le bd$,
the spectral-row part is controlled by the truncated-row estimate,
since $G_i\le C_bd$ and
$a_i>L\sqrt{d/n}$ implies
$\sqrt n\norm{U_{i\cdot}}_2>L/\sqrt{C_b}$.


\subsection{Sharp risk and removal of the logarithmic density floor}

\begin{theorem}[Regularized DC-SBM at every diverging expected-degree
scale]\label{restore:dc-theorem}
Under the fixed structural assumptions \eqref{eq:mean-model}--%
\eqref{eq:structural}, suppose
\[
 d=np\longrightarrow\infty,\qquad np^2\le C_1<\infty.
\]
The regularized algorithm of \cref{restore:algorithm}, using the
DC-SBM profile LRT, satisfies
\[
 \E\ell(\widehat z,z)\le C R_{\mathrm{prof},*}(n)
\]
for all sufficiently large $n$. The constant may depend on the fixed
structural bounds, clamping constants and $C_1$, but not on $n$.
Under the finite-$n$ equal-exposure condition
\eqref{eq:equal-exposure}, the right side may be replaced by
$C R_{\theta,*}(n)$. In particular, the classical balanced planted
partition with fixed separated $q/p$ is included. There is no lower
bound proportional to $\log n$ in this statement.
\end{theorem}

\begin{theorem}[Ordinary SBM with the regularized block-LRT decoder]
\label{restore:sbm-theorem}
Take $\theta_i=1$, fixed arbitrary $K$, comparable group sizes and a
positive compact full-rank block profile as in
\eqref{eq:structural}. Suppose
\[
 d=np\longrightarrow\infty,\qquad np^2\le C_1<\infty.
\]
The ordinary, unprofiled version of \cref{restore:algorithm} satisfies
\[
 \E\ell(\widehat z,z)\le C R_{\mathrm{SBM},*}(n).
\]
This is a sparse block-LRT theorem; it does not replace the broader
dense-probability result for the earlier ordinary decoder, and does not
assert the unchanged free-profile EM schedule in the sparse regime.
\end{theorem}

\begin{proof}[Proof of the sparse-range risk bounds]
Fix any constant $A_0<\infty$ and first work on
$d\le A_0\log n$, still with $d\to\infty$.
Let $\mathcal E$ collect the aggregate geometric conditions of
\cref{prop:sparse-floor}, their recomputed versions after deletion of
at most three vertices, two-sided total degree, and the true block-total event used
in the fitting lemmas. None imposes a simultaneous lower bound on all
raw degrees or a simultaneous incoherence bound. For every prescribed
fixed $r$, its failure probability is $O(n^{-r})$ after increasing the
fixed constants. For the block totals this also follows directly from
the independent-edge Bernstein bound at a sufficiently large constant
times $n$: its exponent is at least $cn/d$, and $np^2\le C_1$ gives
$n/d\ge d/C_1$; in the present range it is much larger than $\log n$.

The initializer and invariant-basin result in
\cref{restore:fit-invariant} apply deterministically on $\mathcal E$.
For a tested target $i$, impose only its upper-degree bound
$D_i\le Ld$, its local normalized row bound, and the condition that its
incident row hits no exceptional deleted eigenvector row. By
\cref{restore:tail-energy-lemma,restore:tail-hit-corollary,restore:tail-target-corollary}, the average probability cost
of the two row conditions is $C_He^{-Hd}$ for any prescribed fixed $H$,
after choosing their fixed thresholds sufficiently large. For the
LOO hit condition, the exceptional set consists of deleted rows with
$\sqrt n\norm{(U_B)_{j\cdot}}_2>R$; the floor implies
\eqref{restore:loo-hit} on its complement. Its expected hit mass
is bounded by the row-count argument in
\eqref{restore:tail-coefficient-count}.
The upper
degree tail has the same property by the single-row Bernoulli Chernoff
bound. Polynomial factors in $d$ are absorbed by choosing a larger
exponential rate at the preceding step. No union over the $n$ tested
targets is taken: their error probabilities are averaged.

On these local events, the fitted-score lemma
\cref{restore:fit-score} replaces every estimated, possibly
data-dependent contrast by a deterministic true block contrast at
$O(1)$ cost. The local LOO comparison
\cref{restore:loo-corrected} replaces its actual full-graph spectral
action by an independent training action at another $O(1)$ cost.
For the DC profile contrast, the new score is the true raw-row
linear profile contrast plus
\[
 \sum_{j\ne i}A_{ij}e_j+O(1),\qquad
 e=(H_B-I)(w-\bar w\one).
\]
The coefficients are measurable from the deleted graph and independent
of the incident Bernoulli row. For each target use the analogous
training event $\mathcal E_i$ formed only from that deleted graph and
its further deletions. In applying the independent-row bound,
drop the full-graph event after the deterministic error inclusion and
condition only on the training graph on $\mathcal E_i$.
The geometry cost is
\[
 \varepsilon_{\rm geom}
 =\PP(\mathcal E^c)+\frac1n\sum_i\PP(\mathcal E_i^c)
 \le C_rn^{-r}.
\]
This does not condition the Bernoulli row on a full-graph event.
The tilted mean and quadratic cost are
bounded by \eqref{eq:floor-tilt-moments}; here $np^2$ is bounded.
The exceptional coefficient hit-mass hypothesis of
\cref{lem:sparse-truncated-transfer} is now a conclusion of
\cref{restore:tail-hit-corollary}, rather than an unproved premise. That sharp
transfer gives, after averaging the local target exceptions and taking
a fixed-$K$ union over competitors,
\[
 \E\ell(\widehat z,z)
 \le \frac Cn\sum_i\sum_{b\ne z_i}
                  d^{-1/2}e^{-I_{i,z_i b}}
       +C_He^{-Hd}+\varepsilon_{\rm geom}.
\]
The same pairwise directional sharp scale is a lower bound, up to a
constant, for the actual profile oracle. Since $K$ is fixed, the
displayed sum is bounded by a constant times its average risk.
Only the averaged target-row tail is used; no pointwise uniform tail
for heterogeneous degree parameters is asserted.

For ordinary SBM, the exact raw Bernoulli contrast has bounded
block-constant log-odds coefficients and a deterministic nonedge
intercept. The fitted intercept error is $O(1)$: each block probability
is estimated at error $O(1/n)$, and the group-size error is $O(n/d)$.
The incident fitted coefficients cost $O(1/d)$ times a reconstructed
row of absolute mass $O(d)$. Replacing wrong labels costs $O(1)$ by
the same weighted wrong-set estimate. Under the exact Bernoulli
Chernoff tilt the probabilities are block constant, so they equal
$D_0h$ with block-constant $\norm h_\infty=O(1/n)$; the tilted
cancellation is exact up to the bounded diagonal-exclusion term.
Thus the same sharp transfer applies to the ordinary true-parameter
oracle, with its nonedge intercept retained. The ordinary fitted and
invariant-basin arguments are also stated explicitly in
\cref{restore:fit-ordinary}.

Uniformly on either compact class, the reference average risk is at
least $c d^{-1/2}e^{-C_Id}$. Choose $H>C_I+1$ and then choose fixed
$r>A_0(C_I+1)+2$. For $d\le A_0\log n$,
\[
 \frac{e^{-Hd}+n^{-r}}{d^{-1/2}e^{-C_Id}}\longrightarrow0.
\]
This closes the unconditional expected-risk estimate throughout the
entire range below any fixed multiple of $\log n$. In particular, no
positive lower coefficient of $\log n$ is required.
\end{proof}

\begin{lemma}[Exponentially accurate aggregate geometry above a logarithmic scale]
\label{restore:high-aggregate}
Fix any $A_0>0$ and $H>0$. Under the fixed positive compact structural
conditions, if $d\ge A_0\log n$ and $np^2\le C_1$, there is an aggregate
event $\mathcal E_H$ with
\[
 \PP(\mathcal E_H^c)\le C_H e^{-Hd}
\]
on which the empirical-floor geometry, weighted degree-energy bounds,
two-sided total-degree bounds, and true block-total bounds used in
\cref{restore:fit-geometry-one,restore:fit-geometry-two,restore:fit-block-event}
hold for the full graph and every graph obtained by deleting at most three
vertices. The constants may depend on $A_0,H,C_1$ and the structural
bounds. No condition is imposed on the smallest raw degree or the
largest spectral row.
\end{lemma}

\begin{proof}
The centered bounded-entry norm inequality of Bandeira--van Handel
\cite[Corollary 3.12 and Remark 3.13]{bvh}, as used in
\eqref{eq:raw-norm-tail}, gives
\[
 \PP\{\norm{A-\Omega}_{\op}>L_H\sqrt d\}
 \le e^{-(H+2)d}
\]
after choosing a sufficiently large fixed $L_H$. Indeed, the row
variance sum is $O(d)$, the entry bound is one, and the prefactor $n$
in the norm-tail inequality is absorbed using $\log n\le d/A_0$.
The completed-mean diagonal has norm $O(p)$ and is absorbed in the
same threshold. This step is solely a norm event; it is not intersected
with the raw minimum-degree event in \eqref{eq:base-event}.

The Bernoulli transform of the total number of edges gives, with
failure probability at most $Ce^{-cnd}$,
\[
 cnd\le\sum_iD_i\le Cnd.
\]
For each of the fixed $K^2$ true block totals, Bernstein's inequality
for the independent upper-triangular edges, whose coefficients have
absolute value at most two, gives
\[
 \PP\left\{\left|[Z^\top(A-\Omega)Z]_{ab}\right|>Ln\right\}
 \le 2\exp\left\{-\frac{cL^2n}{d+L}\right\}.
\]
The deterministic completed diagonal costs $O(d)$ and is included
by increasing $L$. Since $d\to\infty$ and $d^2/n\le C_1$,
the exponent is at least $c'L^2d/C_1$ for all sufficiently large $n$.
Increasing the fixed $L=L_H$ therefore makes the failure of all block
totals at most $C_He^{-(H+2)d}$. The total-degree failure is also
bounded by $Ce^{-(H+2)d}$ for large $n$.

These three full-graph events imply the required events on every
fixed-depth minor without a degree-wise probability union. To see this,
let $J=[n]\setminus T$ with $|T|\le3$, set $A_J=A_{JJ}$ and
$\Omega_J=\Omega_{JJ}$, and recompute all degrees and the empirical
floor on $J$. The centered matrix $A_J-\Omega_J$ is a principal
submatrix of $A-\Omega$, so its norm is at most $L_H\sqrt d$.
Removing $T$ changes the total raw degree by at most $2|T|n$;
as $d\to\infty$, the two-sided bound by fixed multiples of $nd$
therefore persists uniformly over $T$. Removing $T$ changes a true
block total of $A-\Omega$ by $O(|T|n)$ deterministically, since
the entries have bounded absolute value. Its $O_H(n)$ bound persists
as well. The surviving group masses differ from their original
values by $O(1)$ and remain comparable to $n$.

The recomputed average degree is consequently comparable to $d$.
For the fixed sufficiently small floor multiplier $a_0$, its threshold
$\tau_J=a_0\overline D_J$ lies between positive fixed multiples of
$d$ and below every coordinate of $D_{0,J}=\Omega_J\one$.
Thus $G_{J,j}=\max(D_{J,j},\tau_J)\ge cd$ and
$\sum_jG_{J,j}\le Cnd$. The raw norm bound immediately gives
\begin{align*}
 \norm{G_J^{-1/2}(A_J-\Omega_J)G_J^{-1/2}}_{\op}
   &\le C_Hd^{-1/2},\\
 \sum_{j\in J}\frac{(D_{J,j}-D_{0,J,j})^2}{G_{J,j}}
   &=\norm{G_J^{-1/2}(A_J-\Omega_J)\one}_2^2
     \le C_Hn.
\end{align*}
Since the population floor is inactive,
$|G_{J,j}-D_{0,J,j}|\le|D_{J,j}-D_{0,J,j}|$; this proves the
other weighted energy bound. The fixed signal gap and the residual
$\norm{(I-U_JU_J^\top)G_J^{1/2}Z_J}_F\le C_H\sqrt n$
now follow by the deterministic Gram-matrix and projection arguments
in \cref{prop:sparse-floor}. Those arguments use only the floor,
total degree, normalized noise, and degree energy, all just proved.
Intersecting the stated aggregate consequences defines
$\mathcal E_H$. The three full-graph events imply it simultaneously
for every indicated minor, proving its claimed failure bound.

For a deleted graph used as independent training data, form the
analogous event using that graph and its further minors. The same
argument gives failure probability $C_He^{-Hd}$, and this event is
measurable from the training graph alone. The full-graph common event
is used only for deterministic implications, never as a condition
that would preserve independence of an incident row.
\end{proof}

\begin{proof}[Completion over the remaining degree scales]
Fix $A_0>0$ and use the preceding sparse-range argument when
$d\le A_0\log n$. For $d>A_0\log n$, replace its polynomially
accurate aggregate event by $\mathcal E_H$ from
\cref{restore:high-aggregate}. The proof of
\cref{restore:tail-energy-lemma} uses only the deterministic aggregate
geometry on the event; the innovation concentration is applied for
every training graph before restricting to the event. It therefore
holds unchanged with $\mathcal E_H$ and the corresponding
training-measurable events. Its coefficient-hit and target-row
corollaries, the scalar LOO comparison, all fitted-score bounds,
and the invariant basin likewise apply.

The resulting risk upper bound is the appropriate oracle term plus
$C_{H'}e^{-H'd}+C_He^{-Hd}$, where $H'$ can be prescribed by choosing
the fixed local thresholds. On either fixed positive compact model
class the oracle risk is at least $cd^{-1/2}e^{-C_Id}$.
Choosing $H,H'>C_I+1$ pays both remainders. This proves the stated
all-scale bound without a raw information--degree margin or an
entry-ratio-two restriction for the regularized ordinary block-LRT
decoder. The ordinary proof continues to retain its full Bernoulli
nonedge intercept, and the DC proof continues to compare first to
the profile oracle. The known-target-degree DC comparison retains
its separate equal-exposure assumption.
\end{proof}


\subsection{Scope of what has been restored}

The degree-floor row-energy bound, local scalar LOO comparison, and
weighted fitted-parameter argument above repair the three specific
gaps listed at the end of the previous revision's Section 9. The DC
profile-LRT algorithm now has an unconditional constant-oracle-risk
bound at every diverging degree scale in its bounded-$np^2$ class.
Its known-target-degree comparison still requires the stated
equal-exposure restriction; the general obstruction in
\cref{thm:imposs} is unchanged.

For ordinary SBM, the sparse block-LRT route has been restored with an
explicit initializer and true Bernoulli likelihood. This is different
from asserting the previous free-profile soft-EM schedule without
modification. Exact degree preservation and a uniform entry bound
$|R_{ij}|\le Cp$ would force $D_i\le Cd$ for every vertex; that global
bound is not available in the sparse model. The old common-box proof
therefore cannot simply be copied. This algorithm distinction is kept
visible rather than absorbed into a theorem with an unspecified decoder.

All new comparisons are internal mathematical derivations with explicit
proofs. They are not justified by simulations, and they do not claim a
small numerical risk multiplier or a ratio converging to one.
